% !TEX root = ../../../main.tex
% Block card: preprocessing.decompose -- "Decompose" (Part I, the blocks chapter).
\cardfield{What it does}
Splits the trace into three layers that add back up to it exactly: the slow wander
(\emph{baseline}), the fast fuzz (\emph{noise}) and what is left between them (\emph{signal},
where the events are), shows all three, and passes the one you choose on down the chain.

\cardfield{When to use it}
When you want to see what a drift removal and a noise filter would take away before you
commit to them, or when the slow wander or the fuzz is itself what you want to study. It is
the usual first step in front of a detector: pass \emph{signal} on and the detector reads the
events without the drift under them or the fuzz on top. For layers cut by frequency, octave by
octave, use \blockref{preprocessing.wavelet_bands} instead.

\begin{plainterms}[Three layers that add up]
Think of a boat on a tide. The tide raises and lowers everything slowly (the baseline); the
ripples jiggle it many times a second (the noise); the waves you care about sit in between
(the signal). The block measures the tide with a \emph{running median}: at each moment it
takes the middle value of the trace over a long stretch around it, so a short wave does not
move it. It then takes the ripples off what remains with a filter that keeps only fast
changes. Whatever is still left is the signal. Because the signal is defined as ``everything
the other two did not take'', the three always add back up to the recording.
\end{plainterms}

\cardfield{What goes in, what comes out}
In goes a span of one channel with no missing samples. Out comes one layer, as an ordinary
signal of the same length and rate. The block page draws the input and all three layers
stacked on one time axis, under the heading \ui{Every layer}, the one that goes on highlighted.
Click another layer to choose it instead: the click sets \emph{pass\_on}, the same setting the
Parameters panel shows. The row is outlined and the step is marked stale, as after any change
of setting; press \ui{Re-run from} and, once the step has run again, the new layer is
highlighted and the next block reads it.

\cardfield{Settings that matter}
\emph{baseline\_window\_s} is the one to tune first. The running median ignores an event
shorter than about half this window; an event longer than that starts to go into the
baseline instead of the signal. The default 600\,s (ten minutes) suits short events; for drops
that take 5 to 10 minutes to recover, use 1800 to 3600\,s. \emph{noise\_cut\_hz} sets where
fuzz begins: 0, the default, means an eighth of the sampling rate (0.125\,Hz at 1\,Hz, that
is, anything that repeats faster than every 8 seconds). \emph{pass\_on} picks the layer:
\texttt{signal} (the default), \texttt{baseline} or \texttt{noise}.

\cardfield{Watch out for}
A slow event can vanish into the baseline when the window is too short: check the baseline
row before trusting the signal. A sharp edge, such as the start of a sudden drop, also has
fast content, so part of it lands in the noise layer. The first and last half-window of the
span are less reliable. A span with any missing sample is refused, and so is a noise cut at or
above half the sampling rate.
